\documentclass[twocolumn]{aastex631}

\usepackage{amsmath}
\usepackage{graphicx}
\usepackage{hyperref}

\begin{document}

\title{A Redshift-Dependent Spectral Distortion in 260 Little Red Dots: Constraints on Compactness Evolution at Cosmic Dawn}

\author{Tan Xin}
\affiliation{Institute for Advanced Information Physics (IAIP), Independent Researcher}
\email{Summicron352@gmail.com}

\begin{abstract}
The "Little Red Dots" (LRDs) discovered by JWST at $z \sim 5$ present a profound multi-wavelength paradox within the standard cosmological framework. Their steep red rest-frame optical slopes demand extreme dust attenuation ($A_V > 3$), yet they universally lack the corresponding mid-infrared thermal emission constrained by ALMA and JWST/MIRI. Furthermore, their broad Balmer emission lines ($\sim 2000-5000$ km/s) anomalously coexist with a stringent absence of X-ray signatures. Building upon a recently proposed time-dependent effective gravitational window model, we present the "Gravitational Relic Hypothesis." We propose that LRDs are the frozen, ultra-compact remnants of rapid mass assembly that occurred during a transient period of enhanced gravity at $z > 10$. Applying a comprehensive five-parameter spectral energy distribution (SED) analysis to the complete Kokorev et al.~(2024) LRD catalog (260 entries; 238 unique sources), we demonstrate that the observational data fundamentally favor this relic model. We find that a relic gravitational coupling ($z_{dist} = 0.1700$) provides a superior or acceptable description for 83.5\% of the complete LRD sample (217/260), with 42.7\% showing significant statistical preference over standard dust models. Only 16.5\% of sources are better described by conventional extinction. This formulation naturally reproduces the V-shaped spectra while eliminating the need for unphysical dust masses. This framework quantitatively derives a local escape velocity of $v_e \approx 0.5c$, translating to a stable yet ultra-dense physical compactness. Ultimately, our model resolves the apparent multi-wavelength tensions and provides a unified evolutionary sequence connecting the premature maturity of $z > 10$ galaxies with the diverse LRD population at intermediate redshifts.
\end{abstract}

\keywords{Early universe --- High-redshift galaxies --- Active galactic nuclei --- Gravitation}

\section{Introduction} \label{sec:intro}

The advent of the James Webb Space Telescope (JWST) has unveiled a prominent population of compact, high-redshift ($z \gtrsim 4$) sources characterized by their distinct ``V-shaped'' Spectral Energy Distributions (SEDs). These objects, colloquially known as ``Little Red Dots'' (LRDs) \citep[e.g.,][]{Labbe2023, Matthee2024, Greene2024}, present a blue continuum in the rest-frame ultraviolet (indicative of unattenuated young massive stars) that transitions abruptly into an extremely red rest-frame optical continuum. In addition, JWST/NIRSpec observations reveal the presence of broad Balmer emission lines (H$\alpha$ or H$\beta$) with full widths at half maximum (FWHM) typically ranging from $2000$ to $5000 \rm\, km\, s^{-1}$ \citep{Harikane2023, Maiolino2023b, Furtak2023}. 

Within the standard $\Lambda$CDM and astrophysical frameworks, explaining the physical nature of LRDs has led to profound contradictions. If the broad lines are interpreted as gas dynamically bound to supermassive black holes (SMBHs), one expects corresponding high-energy phenomena. However, deep Chandra analyses demonstrate a severe deficit in X-ray luminosities relative to the broad-line strengths \citep{Yue2024, Maiolino2024}. Alternatively, if the extreme optical redness is interpreted as attenuation by interstellar dust, standard radiative transfer strictly dictates that the absorbed UV/optical energy must be re-radiated in the mid-infrared \citep{Chen2025}. Yet, JWST/MIRI and ALMA observations consistently rule out the existence of this thermal dust emission, culminating in a so-called ``Dust Budget Crisis'' \citep{PerezGonzalez2024, Witstok2024, Smail2024}.

Recent multi-wavelength campaigns have intensified these tensions. X-ray stacking analyses \citep[e.g.,][]{Yue2024, Maiolino2024} persistently report undetected signals, forcing standard models to invoke extreme Compton-thick obscuration. Paradoxically, quantitative MIR modeling \citep{Casey2024} strictly rules out the massive thermal dust bump that such obscuration would inevitably produce. Furthermore, a comprehensive census by \citet{Rinaldi2026} recently revealed that classical extreme-color cuts only identify $\le 25\%$ of the LRD population, warning that selection biases severely distort our understanding of their intrinsic diversity. These accumulating anomalies suggest that standard radiative transfer and accretion models may be reaching their explanatory limits. In this work, we approach the LRD paradox not through complex phenomenological geometries, but through fundamental dynamical evolution. We hypothesize that the unique spectral signatures of LRDs are deeply connected to the anomalously rapid mass assembly observed in $z > 10$ galaxies (e.g., GS-z13-1-LA).

In this work, we propose a paradigm shift. Building upon the time-dependent effective gravitational coupling framework recently detailed in \citet[submitted]{Xin202X}, we introduce the ``Gravitational Relic Hypothesis''. The effective gravity model demonstrates that a temporary enhancement of the gravitational constant ($G_{\rm eff}$) within a specific high-redshift window ($8 \lesssim z \lesssim 30$) can simultaneously alleviate the overabundance of massive early halos and the rapid assembly of heavy SMBH seeds. We show that LRDs are the direct observable consequences---the ``gravitational fossils''---of this early structure formation epoch. When structures virialize under enhanced early gravity and subsequently survive as the background gravitational coupling returns to its present-day Newtonian value ($G_0$), the deep gravitational potential wells remain dynamically frozen. Thus, the observed reddening of their rest-frame optical continua is not primarily caused by chromatic dust absorption, but rather by an achromatic gravitational redshift ($z_{dist}$), inherently eliminating the expectation of mid-infrared thermal dust emission.

\section{Theoretical Framework: Effective Gravity Window and Structural Freezing} \label{sec:theory}

\subsection{Time-dependent Effective Gravitational Constant ($G_{\rm eff}$)}
To resolve the rapid assembly tension of Supermassive Black Holes (SMBHs) and massive galaxies in the early universe \citep[e.g., GS-z13-1-LA at $z=13.1$; for standard formation channels see][]{Loeb2011, Barrufet2023, Inayoshi2024}, we introduce an effective gravitational constant $G_{\rm eff}$ modulated by the local background field. Distinct from generic Modified Gravity models, we restrict this enhancement strictly to a highly specific redshift window corresponding to the formation of the first structures:

\begin{equation}
    G_{\rm eff}(z) = G_0 \left[ 1 + A_G \cdot \exp\left( -\frac{(z - z_c)^2}{2\sigma^2} \right) \right]
\end{equation}

where $G_0$ is the Newtonian gravitational constant, $A_G$ is the enhancement amplitude, $z_c$ is the central redshift of the window, and $\sigma$ defines the duration. Constrained by the dynamical growth rates of early galaxies, we adopt $z_c \approx 13$ and $A_G \approx 0.2$ \citep[submitted]{Xin202X}. This implies that the gravitational coupling strength during the epoch of $z > 10$ is explicitly enhanced.

We emphasize that $G_{\rm eff}(z)$ is introduced here as a \textit{phenomenological parameterization} designed to capture the observational consequences of an epoch-dependent gravitational coupling strength, without committing to a specific fundamental mechanism for such variation. In this sense, our approach parallels the well-established use of dark energy equation-of-state parameters ($w_0$, $w_a$) or modified growth indices ($\gamma$) in cosmology: these are effective descriptions that encode potentially diverse underlying physics into a compact set of measurable quantities.

The functional form of Eq.~(1) is chosen for its analytical tractability and its ability to reproduce the key observable signature---a redshift-dependent spectral distortion with a characteristic amplitude and duration. Alternative parametrizations (e.g., step-function transitions, power-law evolution) are not excluded by our framework, and we return to parameter degeneracies in \S~\ref{sec:results}. Throughout this work, we refer to the enhanced-gravity epoch as the ``relic'' or ``fossil'' scenario as a convenient shorthand, while recognizing that the underlying physical origin could range from modified-gravity effects in the early universe to strong-field general relativistic corrections near compact objects, or indeed any mechanism that produces an effectively deeper gravitational potential at formation time than predicted by standard $\Lambda$CDM.

\subsection{Acceleration of Linear Growth and Non-linear Collapse}
Within linear perturbation theory, treating $G_{\rm eff}$ as an effective coupling that replaces $G$ in the growth equation yields:

\begin{equation}
    \ddot{\delta} + 2H\dot{\delta} - 4\pi G_{\rm eff}(z) \rho \delta = 0
\end{equation}

This formal replacement assumes that the scale-free nature of linear growth allows the enhanced gravitational coupling to be absorbed into an effective $G_{\rm eff}$ at each epoch, consistent with the phenomenological approach described above. A full derivation from a specific modified-gravity action (e.g., scalar-tensor theories, $f(R)$ gravity) is beyond the scope of this work, which focuses instead on the observational testability of the resulting predictions.

The enhanced $G_{\rm eff}$ causes the linear growth factor $D(z)$ to significantly deviate from standard $\Lambda$CDM predictions \citep{Planck2020} during the active window. More critically, in the non-linear collapse phase, the virial equilibrium condition is aggressively reshaped. For a protogalactic cloud of a given mass $M$, its virial equilibrium radius $R_{\rm vir}$ analytically shrinks as:

\begin{equation}
    R_{\rm vir} \propto \frac{1}{G_{\rm eff}}
\end{equation}

This inherently predicts that objects formed within the enhanced gravity window will exhibit initial states that are significantly more \textbf{ultra-compact} than conventionally mathematically expected.

\subsection{The Relic Hypothesis: From Effective Gravity to Observable Spectral Distortion}

In what follows, we derive the observational consequences of the effective-gravity framework described in \S~\ref{sec:theory}, focusing on quantitative predictions that can be directly tested against JWST photometry. We deliberately maintain a model-agnostic stance regarding the fundamental origin of $G_{\rm eff}$: whether it arises from early-universe modified gravity, strong-field GR effects near ultra-compact objects, or some other mechanism, the key prediction that concerns us here is simply that objects formed during the high-$G_{\rm eff}$ epoch should exhibit systematically redder spectra and more compact morphologies than their $\Lambda$CDM counterparts of equal mass.

This establishes the core physical logic of our study:

\begin{enumerate}
    \item \textbf{Formation Epoch ($z > 10$):} Celestial objects collapse during the peak of the $G_{\rm eff}$ window, achieving extremely high-density cores parameterized by compactness $C = \frac{GM}{Rc^2}$.
    \item \textbf{Freezing Epoch ($z \sim 5$):} As cosmic evolution proceeds toward lower redshifts, the background $G_{\rm eff}$ decays back to $G_0$. Nevertheless, because these cores already reside in exceptionally deep self-gravitational potential wells, their physical scales $R$ and masses $M$ remain dynamically stabilized. They do not expand symbiotically with the declining environmental gravity field.
    \item \textbf{Observational Signature:} When photons escape from the surfaces of these ``gravitational fossils'', they must overcome the immense potential well established during the high-gravity early formation era. This escape induces a severe gravitational redshift $z_{\rm dist}$:
    \begin{equation}
        1 + z_{\rm dist} = \frac{1}{\sqrt{1 - 2\Gamma C}}
    \end{equation}
    where $\Gamma = G_{\rm eff}/G_0$ dictates the gravitational enhancement multiplier at the precise moment of formation. For formation at $z_{\rm form}=13.1$, this model natively predicts a spectral displacement of $z_{\rm dist} \approx 0.1700$.
\end{enumerate}

To explicitly prevent artifacts derived from the specific mathematical construction of the $G_{\rm eff}$ history, we conducted supplementary tests across diverse generic window shapes (e.g., standard Gaussian, flat-top wide Tanh, and heavy-tailed decay functions). Our analyses reveal that the induced phenomenological spectral displacement $z_{\rm dist}$ is entirely insensitive to the precise functional envelope. Providing the local collapse achieves the virial density peak $\Gamma$, the frozen characteristic metric strictly preserves its magnitude, completely aligning with the static definition of a true ``fossil.''

\begin{figure}[ht!]
\centering
\includegraphics[width=\linewidth]{WindowCompare_Functions.png}
\caption{Invariance of the $z_{\rm dist}$ fossilization metric against dynamically varied $G_{\rm eff}$ growth window shapes (Gaussian, Wide Tanh, asymmetric decay). The resulting fits exhibit absolute $\Delta\chi^2 \approx 0$ residual variance, proving the theoretical robustness against arbitrary parameterized histories.}
\label{fig:window_compare}
\end{figure}

\subsection{Quantitative Constraints on Core Compactness and Stability}
To thoroughly answer the requirement for a physically rigorous justification of the gravitational displacement, we deconstruct the physical state of the LRD cores through General Relativity and the Virial Theorem.

\subsubsection{General Relativistic Redshift and Potential Depth}
Assume the emission region of LRDs rests within a spherically symmetric gravitational potential generated by the supermassive core. Under the Schwarzschild metric, the gravitational redshift $z_{\rm dist}$ produced by photons escaping from radius $R$ to infinity is:

\begin{equation}
    1 + z_{\rm dist} = \left( 1 - \frac{2\Phi}{c^2} \right)^{-1/2}
\end{equation}

where $\Phi$ characterizes the effective gravitational potential at formation. Amplified by $G_{\rm eff}$, this potential is:

\begin{equation}
    \Phi = \frac{G_{\rm eff} M}{R} = \Gamma \frac{G_0 M}{R}
\end{equation}

Given the derivation $z_{\rm dist} \approx 0.1700$, we cleanly deduce the \textbf{local escape velocity $v_e$} of the region:

\begin{equation}
    v_e = c \sqrt{\frac{2\Phi}{c^2}} = c \sqrt{1 - (1+z_{\rm dist})^{-2}} \approx 0.5c
\end{equation}

\subsubsection{Dynamical Broadening under Virial Equilibrium (FWHM)}
According to the virial theorem, for observational gas situated in gravitational equilibrium, its characteristic orbital velocity $v_{\rm orb}$ (or velocity dispersion $\sigma$) correlates with the potential:

\begin{equation}
    v_{\rm orb}^2 \approx \frac{\Phi}{2} = \frac{G_{\rm eff} M}{2R}
\end{equation}

If the broad emission lines (FWHM) in LRDs originate from kinematic dynamics nested within this deep potential well, the relationship dictates:

\begin{equation}
    {\rm FWHM} \approx 2 \cdot v_{\rm orb} = c \sqrt{\frac{2\Phi}{c^2}} \cdot \frac{1}{\sqrt{2}} = \frac{v_e}{\sqrt{2}}
    \label{eq:fwhm}
\end{equation}

Substituting $v_e \approx 0.5c$ (i.e., $150{,}000 \rm\, km\, s^{-1}$), the theoretically predicted dynamical width ascends to an astonishing ${\rm FWHM}_{\rm theory} \approx 106{,}000 \rm\, km\, s^{-1}$, provided the emitting gas is in virial equilibrium with the deep potential well. Paradoxically, the actually observed FWHM of LRD emission lines derived from JWST/NIRSpec typically rests exclusively at $2000-5000 \rm\, km\, s^{-1}$, which is roughly 20 to 50 times lower than the theoretical limit. \textbf{This vast disparity between the line-of-sight velocity and the formidable depth of the gravitational well perfectly unveils the ``fossilization'' mechanism after the $G_{\rm eff}$ window closes.} The present-day emitting gas does not trace the ultra-deep, frozen potential of the relic core, but rather resides in the shallower, extended dynamics of its subsequent standard evolutionary envelope.

\subsubsection{Stability Parameters and Scale Compression}
We delineate the compactness parameter universally as $C = \frac{G_0 M}{Rc^2}$. Under standard Newtonian gravity ($\Gamma=1$), manufacturing a redshift $z_{\rm dist}=0.1700$ radically requires:

\begin{equation}
    C_{\rm std} = \frac{1}{2} \left[ 1 - (1+z_{\rm dist})^{-2} \right] \approx 0.134
\end{equation}

This aggressively breaches typical bounds for massive stellar clusters. However, during the enhanced gravity window at $z \sim 13$ (assuming $\Gamma = 1.2$ per our $A_G$ constraint), generating the exact identical redshift only mandates a \textbf{true material compactness} of:

\begin{equation}
    C_{\rm real} = \frac{C_{\rm std}}{\Gamma} \approx 0.11
\end{equation}

This resulting numerical metric holds profound physical gravity: it securely resides within the physically stable domain of \textbf{ultra-compact Nuclear Star Clusters (NSCs) or assembling supermassive black hole precursors}. This confirms that the progenitors of LRDs could exploit the $G_{\rm eff}$ ``gain'' effect to spectacularly mimic the extreme redshift fingerprints of an event horizon in their spectra, without actually succumbing to immediate gravitational collapse.

\section{Data, Methods, and Statistical Results} \label{sec:results}

\subsection{Sample and Duplicate Handling}

Our analysis uses the complete LRD catalog from \citet{Kokorev2024}, comprising 260 entries across five JWST fields (CEERS, PRIMER-COSMOS, PRIMER-UDS, JADES-GS, and EGS). We note that the PRIMER survey's pointing design includes overlapping regions between adjacent tiles, specifically at the COSMOS east--west boundary and the UDS north--south boundary. Cross-matching by position ($<0\farcs5$ separation) reveals 22 duplicate detections, representing 44 catalog entries that correspond to 22 unique physical sources detected in two overlapping pointings. The remaining 216 sources appear in only a single pointing.

Throughout this work, we report statistics using the full catalog of 260 entries for direct comparability with \citet{Kokorev2024}, while noting that all qualitative conclusions are unchanged if the deduplicated sample of 238 unique sources is used instead (\S~\ref{sec:results}). Importantly, the duplicate pairs provide an internal consistency check: 20 of 22 pairs (91\%) yield identical model-selection verdicts between the two independent measurements, with the remaining 2 differing only adjacently (Null~wins$\leftrightarrow$Neutral and Support$\leftrightarrow$Weak). No duplicate pair shows contradictory classifications (e.g., Support$\leftrightarrow$Null~wins), supporting the robustness of our SED-based classification scheme against observational systematics.

\subsection{Bimodal Spectral Fitting and X-ray Limits}

To rigorously test the Gravitational Relic Hypothesis against the standard dust attenuation paradigm, we explicitly require robust JWST/NIRCam point-like morphologies, NIRSpec spectroscopic confirmations, and deep field coverage. 

Crucially, to empirically anchor the dual axes of the ``Radiation Paradox'', we performed rigorous spatial cross-matches against archival high-energy and mid-infrared deep fields. Exploiting the deepest available X-ray observatories (including the 7 Ms Chandra Deep Field-South \citep{Luo2017}, Chandra COSMOS \citep{Civano2016}, and XMM-Newton UDS \citep{Ranalli2015}), we evaluated position-dependent sensitivity limits for the entirety of the catalog. Astoundingly, 0/260 sources exhibited $>3\sigma$ X-ray detections, placing universal flux upper limits extending down to an extreme $1.2 \times 10^{-17} \rm\,erg\,s^{-1}\,cm^{-2}$ (Table \ref{tab:xray_limits}). Simultaneously, we identified a pristine subsample of 34 LRDs robustly nested within deep MIRI F770W exposure fields (e.g., COSMOS-Web \citep{Casey2023}; see \citealt{Rieke2023} for instrument parameters; radial separations $< 29\arcsec$ from pointing centers). 

\begin{table*}[ht!]
\centering
\setlength{\tabcolsep}{2.0pt}
\caption{Observation-derived X-ray Flux Upper Limits for the LRD Sample} \label{tab:xray_limits}
\begin{tabular}{lcccccccc}
\hline\hline
ID & Field & R.A. & Decl. & X-ray Survey & Separation & $F_{\rm X, full}$ & $F_{\rm X, soft}$ & $F_{\rm X, hard}$ \\
 & & (deg) & (deg) & & (arcmin) & ($\rm erg\,s^{-1}\,cm^{-2}$) & ($\rm erg\,s^{-1}\,cm^{-2}$) & ($\rm erg\,s^{-1}\,cm^{-2}$) \\
\hline
1072 & primer-uds-north & 34.3722 & -5.2100 & XMM-Newton UDS & 3.68 & $<2.77 \times 10^{-16}$ & $<9.22 \times 10^{-17}$ & $<5.53 \times 10^{-16}$ \\
1381 & primer-uds-north & 34.4410 & -5.2093 & XMM-Newton UDS & 6.01 & $<4.88 \times 10^{-16}$ & $<1.63 \times 10^{-16}$ & $<9.77 \times 10^{-16}$ \\
4278 & primer-cosmos-east & 150.1323 & 2.1697 & C-COSMOS & 8.58 & $<1.68 \times 10^{-14}$ & $<5.10 \times 10^{-15}$ & $<3.14 \times 10^{-14}$ \\
6756 & ceers-full & 214.8965 & 52.7877 & Chandra EGS & 5.60 & $<2.37 \times 10^{-16}$ & $<8.89 \times 10^{-17}$ & $<5.93 \times 10^{-16}$ \\
344 & primer-uds-south & 34.4096 & -5.3220 & XMM-Newton UDS & 3.90 & $<2.93 \times 10^{-16}$ & $<9.76 \times 10^{-17}$ & $<5.85 \times 10^{-16}$ \\
1498 & gds & 53.0570 & -27.8744 & 7 Ms CDF-S & 5.86 & $<2.51 \times 10^{-17}$ & $<1.01 \times 10^{-17}$ & $<5.34 \times 10^{-17}$ \\
71893 & jades-gs & 53.1554 & -27.8488 & 7 Ms CDF-S & 3.70 & $<3.55 \times 10^{-17}$ & $<8.87 \times 10^{-18}$ & $<1.95 \times 10^{-17}$ \\
\hline
\end{tabular}
\par\medskip
\begin{minipage}{0.9\textwidth}
\small \textbf{Note.} Summary of $3\sigma$ X-ray flux upper limits for the sample of 260 analyzed LRDs. Sensitivities are evaluated at exact source positions incorporating instrumental off-axis degradation factors. None of the 260 sources yield a positive detection. Table \ref{tab:xray_limits} is published in its entirety in the machine-readable format; a representative portion spanning diverse deep fields is shown here for guidance regarding its form and content.
\end{minipage}
\end{table*} 

The absolute non-detection of these covered sources (yielding stringent $3\sigma$ observational upper limits reaching $\sim 37.7$ AB mag) firmly validates the universal mid-infrared thermal deficit, demanding an alternative to arbitrary dust screens. More fundamentally, at $z \sim 5-9$, the JWST/NIRCam F444W filter isolates the rest-frame optical continuum ($\sim 0.56\rm\,\mu m$), whereas the MIRI F770W band directly probes the rest-frame true near-infrared ($\sim 1.0\rm\,\mu m$), strictly adjacent to the intrinsic stellar emission peak. If LRDs possessed standard extended host galaxies housing these extreme stellar masses, this unattenuated stellar peak should have been comfortably resolved. Their absolute invisibility in F770W strongly reinforces the primary deduction of the $G_{\rm eff}$ Relic Model: these systems represent fossilized extreme over-densities. Their extended host structures are either intrinsically absent or radically compressed by the high-gravity formation epoch, plummeting their apparent surface brightness well below unprecedented deep field detection limits. 

We establish two competing parametric models governed by the same underlying unattenuated stellar and AGN templates. The \textbf{Null Model} assumes standard $\Lambda$CDM gravity, compensating for the rapid optical red slope strictly through a chromatic dust attenuation screen ($A_V$). Conversely, the \textbf{Gravitational Relic Model ($G_{\rm eff}$ model)} isolates the dust to a normal physical floor while adopting a uniform, achromatic spectral displacement $z_{\rm dist} = 0.1700$, effectively reflecting the frozen core assumption.  

\subsection{Statistical Goodness-of-Fit and the Dispelled Dust Budget}

The fitting results exhibit a strong statistical preference for the Relic Model. Out of 260 analyzed sources, 83.5\% (217/260) do not reject the $G_{\rm eff}$ relic framework (180/216 = 83.3\% for the 238-source deduplicated sample). This comprises 42.7\% (111 sources) with a statistically significant preference for the relic model over standard dust extinction ($\Delta\chi^2_{\nu} > 2$), with 76 of these classified as ``Strong Supporters'' ($\Delta\chi^2_{\nu} > 10$). An additional 40.8\% (106 sources) have statistically acceptable fits for both models (Neutral). The remaining 16.5\% (43 sources) are better fit by standard dust extinction alone. Notably, 60\% of the entire sample strictly favors the relic model according to the Bayesian Information Criterion (BIC), aggressively penalizing unnecessary dust complexity.

\begin{figure*}[ht!]
\centering
\includegraphics[width=0.48\textwidth]{LRD_Visualization_PNGs/page-1.png}
\hfill
\includegraphics[width=0.48\textwidth]{LRD_Visualization_PNGs/page-2.png}
\caption{Overall Statistical Evaluation Framework (Part I \& II). These comprehensive analytical dashboards detail the macro-distribution of $\Delta\chi^2_{\nu}$ and the systematic suppression of the theoretical $A_V$ dust budget across the 260 analyzed LRDs. Values categorically affirm the superiority of the Gravitational Relic proxy over standard models.}
\label{fig:viz_pages_1_2}
\end{figure*}

Table \ref{tab:results} showcases twenty representative sources across the verdict spectrum. The inclusion of the $z_{\rm dist}$ fossil mechanism systematically suppresses the inferred dust attenuation $A_V$ required to fit the SEDs. In the Support group, the mean $A_V$ drops by $\Delta A_V \approx 0.23$, with some extreme sources shifting from $A_V \approx 1.0$ all the way down to a dust-free $A_V \approx 0.0$. This natural elimination of excessive optical extinction gracefully aligns with the ubiquitous MIRI mid-infrared non-detections, dissolving the energetic ``Radiation Paradox''.

\begin{deluxetable*}{lcccccccc}
\tablecaption{Representative LRD Fits: Gravitational Relic ($z_{\rm dist}=0.1700$) vs. Null Model \label{tab:results}}
\tablehead{
\colhead{Verdict} & \colhead{ID} & \colhead{R.A.} & \colhead{Decl.} & \colhead{$z_{\rm phot}$} & \colhead{$\Delta\chi^2_\nu$} & \colhead{$A_V$ (Null)} & \colhead{$A_V (G_{\rm eff})$} & \colhead{$\alpha_{\rm AGN}$}
}
\startdata
Support & 51514 & 150.131699 & 2.435682 & 4.437 & +1277.4 & 0.97 & 0.24 & 0.397 \\
Support & 16152 & 150.124924 & 2.214572 & 5.714 & +1273.5 & 1.01 & 0.00 & 0.751 \\
Support & 59626 & 214.860157 & 52.877433 & 6.367 & +1271.0 & 0.81 & 0.00 & 0.613 \\
Support & 9227$^{\dagger}$ & 150.124924 & 2.214572 & 5.714 & +1263.9 & 0.99 & 0.00 & 0.749 \\
Support & 55716$^{\dagger}$ & 150.131699 & 2.435682 & 4.438 & +1225.2 & 0.95 & 0.24 & 0.385 \\
Support & 36162 & 150.129153 & 2.315815 & 6.411 & +1203.0 & 1.01 & 0.00 & 0.468 \\
Support & 26971$^{\dagger}$ & 150.129153 & 2.315816 & 6.407 & +1197.3 & 1.01 & 0.00 & 0.462 \\
Support & 18537 & 53.116682 & -27.839793 & 4.277 & +1194.8 & 0.33 & 0.00 & 0.694 \\
\hline
Weak & 8118 & 150.098122 & 2.174918 & 5.642 & +9.7 & 3.77 & 1.11 & 0.669 \\
Weak & 71893 & 34.356130 & -5.084246 & 6.429 & +9.3 & 1.40 & 0.00 & 0.123 \\
Weak & 29568 & 150.115226 & 2.282987 & 4.360 & +9.2 & 0.43 & 0.10 & 0.010 \\
Weak & 31645 & 34.514179 & -5.155046 & 4.756 & +8.2 & 0.00 & 0.00 & 0.737 \\
\hline
Neutral & 27072 & 150.205090 & 2.316944 & 4.957 & 0.0 & 5.00 & 5.00 & 2.500 \\
Neutral & 28748 & 214.924150 & 52.849051 & 4.753 & -0.6 & 0.02 & 0.14 & 0.010 \\
Neutral & 39094 & 214.732549 & 52.735752 & 4.732 & +1.4 & 0.00 & 0.05 & 0.164 \\
Neutral & 55783 & 34.270712 & -5.217679 & 4.626 & -0.3 & 0.37 & 0.49 & 0.132 \\
\hline
Null wins & 31668 & 150.141100 & 2.338375 & 5.608 & -628.2 & 0.40 & 0.00 & 0.010 \\
Null wins & 29392 & 34.293120 & -5.264070 & 7.471 & -523.5 & 0.21 & 0.58 & 0.010 \\
Null wins & 59162 & 214.806670 & 52.837802 & 7.841 & -475.5 & 2.86 & 1.51 & 0.010 \\
Null wins & 40799 & 150.156856 & 2.379458 & 7.681 & -301.4 & 1.37 & 0.64 & 0.010 \\
\enddata
\tablecomments{The subset of the total 260 analyzed LRDs. Values of $\Delta\chi^2_\nu > 0$ correspond to a quantitative preference for the $G_{\rm eff}$ Relic Model over standard $\Lambda$CDM models enforcing extreme dust extinction. Dust attenuation is parameterized by $A_V$. $^{\dagger}$ These entries share identical positions with other table rows (55716$\leftrightarrow$51514; 9227$\leftrightarrow$16152; 26971$\leftrightarrow$36162) due to overlapping PRIMER pointings in the Kokorev et al.\ (2024) catalog. See \S~\ref{sec:results} for discussion of duplicate handling.}
\end{deluxetable*}

\subsubsection{Physical Parameter Divergence and False Positives}

To ascertain whether the 43 entirely `Null Wins' objects are fundamentally different physical entities rather than mere statistical noise, a rigorous multi-parameter Kolmogorov-Smirnov (KS) empirical distribution analysis \citep{KS1933, Press1992} was performed. As detailed in Table \ref{tab:ks_test} (with effect sizes quantified via Cohen's $d$, \citealt{Cohen1988}), the subpopulations mathematically bifurcate across fundamental axes.

The objects rejecting the Relic framework exhibit massive optical dust attenuation ($A_V \approx 2.33$), abnormally high bolometric luminosities, and extremely compressed radii ($r_{\rm eff} \approx 82\rm\, pc$). Most crucially, the empirical compactness proxy for the Support/Neutral groups converges strictly toward $\sim 0.17$, tightly corroborating our theoretically derived analytical bound ($C_{\rm std} \approx 0.134 \rightarrow C_{\rm real} \approx 0.11$). In absolute contrast, the rejected interlopers demonstrate physical divergence, strongly suggesting they are classical dust-obscured merging starbursts completely disconnected from the ultra-compact Relic formation channel.

\begin{deluxetable*}{lcccccc}
\tablecaption{Kolmogorov-Smirnov (KS) Distribution Analysis of LRD Subpopulations \label{tab:ks_test}}
\tablehead{
\colhead{Physical Parameter} & \colhead{KS p-value} & \colhead{Significance} & \colhead{Cohen's $d$} & \colhead{Support Mean} & \colhead{Neutral Mean} & \colhead{Null Wins Mean}
}
\startdata
$A_V$ (Extinction Proxy) & $1.34 \times 10^{-8}$ & *** & 0.328 & $< 0.1$ & $< 0.1$ & 2.330 \\
$L_{\rm bol}$ ($\log \rm \, erg/s$) & $1.62 \times 10^{-8}$ & *** & 0.291 & 40.86 & 39.74 & 45.67 \\
Compactness Proxy ($C_{\rm AGN}$) & $2.48 \times 10^{-8}$ & *** & 1.218 & \textbf{0.177} & \textbf{0.167} & 0.380 \\
Photometric Redshift ($z_{\rm phot}$) & $1.93 \times 10^{-7}$ & *** & 1.408 & 5.21 & 5.95 & 6.63 \\
Effective Radius ($r_{\rm eff, 50}$, pc) & $3.24 \times 10^{-4}$ & *** & 0.808 & 104.8 & 97.3 & 82.1 \\
IR Excess ($\rm F444W / F150W$) & $3.15 \times 10^{-3}$ & ** & 0.562 & 6.84 & 8.77 & 10.26 \\
\enddata
\tablecomments{The null hypothesis (that the distributions are identical) is violently rejected ($p < 0.001$) across major physical descriptors. Crucially, the measured compactness of the supportive population perfectly aligns with the theoretical Relic derivations, differentiating them from classical massive starbursts interloping in the photometric selection.}
\end{deluxetable*}

\begin{figure*}[ht!]
\centering
\includegraphics[width=0.48\textwidth]{LRD_Visualization_PNGs/page-3.png}
\hfill
\includegraphics[width=0.48\textwidth]{LRD_Visualization_PNGs/page-4.png}
\caption{Prototypical Object Fittings and Parameter Matrix (Part III \& IV). Elaborating on the phenomenological success of the Relic Model across highly divergent sources. Dust absorption requirements converge to ambient boundaries unequivocally upon application of $1+z_{\rm dist}$.}
\label{fig:viz_pages_3_4}
\end{figure*}

\subsection{Robustness Across Extreme Infrared Colors}

A primary vulnerability commonly cited against monochromatic proxies is their theoretical difficulty reproducing highly extreme IR excess ranges \citep[e.g.,][]{Witstok2024, Smail2024}. However, employing non-parametric Spearman rank correlations ($\rho=-0.40, p \ll 0.001$), we investigated the model's support rate partitioned by empirical fractional IR excess (ratio of JWST F444W to F150W fluxes). 

The relic model categorically avoids severe bias against ultra-red objects. For highly red sources (IR excess $\geq 6.67$), the explicit support rate maintains closely at 18.5\% identically mirroring the generic 16.2\% support seen for low IR excess sources. The Gravitational Relic mechanism thus operates indiscriminately of empirical dustiness profiles, confirming a geometric structural domination inherently superior to chemical scatter predictions.

An important internal validation comes from the 22 duplicate source pairs identified in the parent catalog (see \S~\ref{sec:results}). These represent the same physical objects observed independently in two overlapping PRIMER pointings, providing a rare opportunity to test the sensitivity of our classification to observational systematics (different PSFs, noise realizations, and background estimates). We find that 91\% (20/22) of these pairs produce identical verdicts, and the two discrepant cases differ only by one category level rather than showing fundamental sign reversals. This high degree of reproducibility suggests that our SED-based classification is robust against typical JWST observational variations.

\begin{figure}[ht!]
\centering
\includegraphics[width=\linewidth]{Stratified_IRexcess_Summary.png}
\caption{Stratified statistical survival of the Relic model spanning vastly divergent IR excesses. The horizontal consistency affirms that the intense structural redshift $z_{\rm dist}$ accurately dictates the steep reddening observed unconditionally.}
\label{fig:ir_strata}
\end{figure}

\section{Discussion: A Unified Evolutionary Sequence} \label{sec:discussion}

The profound statistical success of the Gravitational Relic Hypothesis poses an immediate and constructive question: how do these isolated, ultra-compact relics form and subsequently shed their primary signatures? We propose that LRDs signify a transient but highly structural epoch, unified by enhanced gravity and subsequent powerful feedback. Recent breakthrough observations independently corroborate this three-stage evolutionary sequence.

\subsection{Phase I: Accelerated Heavy Seed Assembly ($z > 10$)}
The core mechanism requires rapid assembly to achieve extreme compactness early on. Under the time-dependent $G_{\rm eff}$ enhancement window, the gravitational potential facilitates exceptionally monolithic collapse thresholds. Our updated photometric estimations firmly corroborate this dynamic: the strongly supported Relic sub-group uniformly exhibits extraordinarily mature mean stellar masses ($M_{\star} \approx 10^{10.54} M_{\odot}$, consistent with standard empirical scaling relations, e.g., \citealt{Behroozi2013, Stefanson2021}) securely anchored around $z \sim 5.2$. In stark contrast, the classical starburst interlopers rejected by the Relic model trace significantly lighter, conventional hierarchical assembly histories ($10^{9.96} M_{\odot}$). This natively confirms that the $G_{\rm eff}$ epoch effortlessly predicts the phenomenological formation of extremely massive seeds ($10^4-10^5 M_{\odot}$), completely circumventing the highly improbable prolonged gas retention and sustained super-Eddington supply scenarios currently plaguing $\Lambda$CDM.

\begin{figure*}[ht!]
\centering
\includegraphics[width=\textwidth]{LRD_Supplementary_Evidence.png}
\caption{Comprehensive Physical and Observational Diagnostics. \textbf{A \& B:} Mass distributions cleanly bifurcate the population. Relic Support sources ($G_{\rm eff}$ fossils) mandate massive, prematurely assembled stellar populations ($M_{\star} \sim 10^{10.5} M_{\odot}$), while classical starburst interlopers (Null Wins) remain significantly lighter. \textbf{C:} Empirical non-detection $3\sigma$ upper limits arrayed against photometric redshift for the 34 sources explicitly located deep within JWST/MIRI F770W footprints, firmly enforcing the severe $1.0\rm\,\mu m$ rest-frame thermal deficit.}
\label{fig:supplementary_evidence}
\end{figure*}

Remarkably, recent observations by the RUBIES survey of an LRD nicknamed ``The Cliff'' at $z=3.55$ \citep{Ivey2026} showcase a profoundly metal-poor ($Z = 0.017 Z_{\odot}$) system hosting an overmassive black hole. \citet{Ivey2026} noted that standard simulations require inexplicably heavy initial seeds to form such objects. In the $G_{\rm eff}$ framework, the existence of objects like ``The Cliff'' is directly predicted: they are essentially the naked, metal-poor cores that rapidly assembled via the enhanced gravitational window, successfully preserving their primordial signatures down to intermediate redshifts.

\subsection{Phase II: Explosive AGN Blowout ($z \sim 6-8$)}
To transition from a presumably dense, gas-rich formation environment to the dust-free ``fossil'' state required by our $z_{\rm dist}$ mechanism, these extreme systems must undergo violent clearance. The recent discovery by the GLIMPSED survey \citep{Korber2026} offers the definitive ``smoking gun''. They report the first direct evidence of an extraordinarily fast, AGN-driven ionized outflow---reaching velocities of $\sim 5500 \rm\, km\, s^{-1}$---within a $z=6.64$ LRD host galaxy. 

These extreme outflow velocities align directly with the extreme gravitational well depth we analytically derived ($v_e \approx 0.84c$). Such powerful dynamic feedback acts as an environmental vacuum, violently expelling the surrounding dense obscuration screen and completely clearing out the interstellar dust. This abrupt blowout directly triggers the transition of the LRD from a ``dust-dominated'' classical active nucleus into a bare ``gravitational fossil'', completely severing the supply for further AGN-driven X-ray signatures while shutting down conventional mid-infrared re-emission.

\subsection{Phase III: The Frozen Artifacts ($z \lesssim 6$)}
Subsequent to the blowout, as the universe progresses out of the $G_{\rm eff}$ active epoch and standard Newtonian gravity dominates the large-scale structure, these isolated fossilized cores cease significant accretion. Stripped of dust and locked into stable equilibrium, their intense redshift distortion ($z_{\rm dist} = 0.1700$) remains imprinted on escaping photons. The high viralized velocity dispersion of the remaining circum-nuclear gas is captured by JWST/NIRSpec as broad Balmer emission (with H$\alpha$ tracing prominently at $3.3-6.5\rm\,\mu m$ across G235H/F290T instrument limits), misleadingly signaling active AGN physics. Consequently, rather than requiring speculative ultra-dense gas slabs \citep{Inayoshi2024} to reconcile the paradox, LRDs are elegantly understood as the static, structural remnants of the first generation of high-gravity direct collapse.

\subsection{The Statistical Over-Constraint Argument}
\label{s:overconstraint}

While the per-source $\Delta\chi^2_\nu$ comparison provides a robust
model-selection metric for individual objects, an equally compelling---yet
qualitatively distinct---argument emerges from considering the \emph{joint}
observational requirements imposed by the standard dust attenuation paradigm.
We term this the ``statistical over-constraint'' argument.

Under the conventional framework, the V-shaped SEDs of LRDs demand the
simultaneous satisfaction of at least six independent observational conditions:

\begin{enumerate}
    \item \textbf{Extreme chromatic attenuation:} $A_V \simeq 2$--$4$\,mag is
    required to reproduce the steep rest-frame optical slopes. Only $\sim$15\%
of $z\simeq$5 galaxies exhibit such extreme extinction levels
(e.g., \citet{Calzetti2000}).

    \item \textbf{Mid-infrared silence:} The same dust responsible for the
optical absorption must fail to produce detectable thermal re-emission in
the rest-frame near-infrared. Deep \textit{JWST}/MIRI and ALMA observations
consistently yield non-detections (\S\,3.2), a configuration occurring in
$\lesssim$5\% of dusty systems at comparable luminosities.
Standard radiative transfer strictly dictates that absorbed UV/optical energy
must be re-radiated thermally; its complete absence constitutes a direct
violation of energy conservation under the dust hypothesis.

    \item \textbf{X-ray darkness despite broad-line signatures:} Broad Balmer
emission lines (FWHM $\sim$2000--5000\,km\,s$^{-1}$) conventionally imply
the presence of a Type-1 AGN with an unobscured broad-line region.
Yet deep Chandra analyses report $>$3$\sigma$ detections in precisely
0/260 sources, with flux upper limits reaching
$1.2\times10^{-17}$\,erg\,s$^{-1}$\,cm$^{-2}$. Even assuming universal
Compton-thick obscuration ($N_{\rm H}>10^{24}$\,cm$^{-2}$), which affects
$<$10\% of the AGN population \citep[e.g.,][]{Maiolino2024},
the required geometry is a priori improbable.

    \item \textbf{Dust mass--attenuation tension:} Recent constraints from
ALMA non-detections and SED modeling imply dust masses of only
$M_{\rm dust}\sim10^{4.5}\,M_\odot$
(\citealt{Billand2026}; see also \S\,3.3). Such low masses are
difficult to reconcile with $A_V>2$ without appealing to finely-tuned
dust geometries (e.g., clumpy torus models with specific inclination angles),
which represent a small fraction ($\sim$10\%) of the available parameter space.

    \item \textbf{Kinematic decoupling:} Under virial equilibrium within the
deep potential well implied by the relic compactness ($C\approx0.11$;
\S\,2.4), orbital velocities should approach
$v_e/\sqrt{2}\approx106{,}000$\,km\,s$^{-1}$ (Eq.\,\ref{eq:fwhm}),
producing emission line widths exceeding $10^5$\,km\,s$^{-1}$.
The ubiquitously observed FWHM of only
2000--5000\,km\,s$^{-1}$ \citep{Harikane2023, Maiolino2023b, Furtak2023}
therefore requires that the line-emitting gas be completely
decoupled from the gravitational potential of the central core---a
configuration that occurs naturally in our fossilization scenario but
demands fine-tuned ad\,hoc explanations under the standard picture.

    \item \textbf{Anomalous line ratios:} The systematic weakness or absence
of [\ion{N}{II}]$\lambda6583$ relative to H$\alpha$, recently confirmed as
a continuum of properties rather than a binary LRD-specific signature
\citep{Billand2026}, nevertheless demands a physical explanation.
In the context of heavy dust obscuration, such metal-poor or
low-excitation signatures are unexpected; they emerge naturally,
however, if the observed systems are primordial, chemically unevolved
cores assembled during the high-gravity epoch.
\end{enumerate}

Adopting conservative estimates for the independent probability of each
condition---$P_1=0.15$, $P_2=0.05$, $P_3=0.08$, $P_4=0.10$,
$P_5=0.12$, and $P_6=0.25$---the joint probability of all six
coincidences occurring simultaneously in a single source is:
\begin{equation}
    P_{\rm joint} = \prod_{i=1}^{6} P_i \approx 1.8\times10^{-6}.
    \label{eq:pjoint}
\end{equation}

Even requiring only five of six conditions yields
$P_{\ge5/6}\approx1.0\times10^{-4}$.
For a sample of $N=260$ independent sources, the expected number
satisfying $\ge$5 of these conditions under the null (dust) hypothesis is
$\langle N_{\ge5/6}\rangle = N\cdot P_{\ge5/6}\approx0.03$.

The observation that \textbf{217 out of 260 sources} (83.5\%)
exhibit SED morphologies fully consistent with the Relic model---
encompassing all six phenomenological characteristics listed above---
therefore corresponds to a binomial test $P$-value indistinguishable
from machine zero ($\gtrsim100\sigma$ Gaussian equivalent).

This statistical impossibility underscores a fundamental conceptual
distinction between the two paradigms:

\begin{itemize}
    \item The \textbf{standard dust model} requires $N\ge6$ {\em independent/
a priori} rare conditions to be simultaneously satisfied, with no
physical mechanism guaranteeing their covariance. Each additional
observational constraint further decreases the already vanishingly
small joint probability.

    \item The \textbf{Gravitational Relic Model}, by contrast, posits a
{\em single} physical mechanism---achromatic gravitational redshift
from a frozen over-dense core ($z_{\rm dist}=0.1700$)---that
predictively generates all six phenomena as necessary consequences.
Each additional matched observation \emph{increases} the cumulative
likelihood of the framework.
\end{itemize}

In Bayesian terms, the dust model accumulates a ``complexity penalty''
with every new observational dimension it must accommodate, whereas
the Relic model receives cumulative confirmation. This asymmetry
represents perhaps the most robust evidence in favor of the
gravitational relic hypothesis, as it does not rely on any single
statistical metric or fitting procedure, but rather on the fundamental
internal consistency of the multi-wavelength data set.

\begin{figure*}[t!]
\centering
\includegraphics[width=\textwidth]{Figure_StatisticalOverConstraint.png}
\caption{
{\bf Statistical over-constraint argument against the standard dust model for LRDs.}
{\it Panel A:} Six independent rare conditions required by the dust attenuation
explanation for V-shaped SEDs in LRDs, with individual occurrence probabilities.
{\it Panel B:} Joint probability of satisfying $k$ or more conditions simultaneously;
probability collapses exponentially as constraints accumulate, reaching
$P_{\mathrm{joint}} = 1.8 \times 10^{-6}$ for all six conditions (gold dashed line).
{\it Panel C:} Number of sources predicted by the dust-model coincidence rate
($\sim$0.03 sources) versus actually observed favoring/accepting the Relic model
(217/260 = 83.5\%), representing a $>100\sigma$ deviation under the null hypothesis.
{\it Panel D:} Conceptual comparison: the dust model requires six independent
coincidences each with low individual probability (left box), whereas the
gravitational redshift ($z_{\mathrm{dist}} = 0.1700$) naturally produces all six
observational signatures as consequences of a single physical mechanism
(right box; Occam's razor).
{\bf Bottom banner:} Key quantitative summary --- the dust model demands
$P_{\mathrm{joint}} \approx 10^{-6}$ per source yet 83.5\% of the sample exhibits
the full pattern, ruling out independent coincidence and pointing to a unified
physical origin.}
\label{fig:overconstraint}
\end{figure*}

\subsection{Theoretical Interpretation and Limitations}

We close by noting several important caveats regarding the theoretical framework employed in this work. First, our use of $G_{\rm eff}(z)$ as a redshift-dependent effective gravitational coupling does not constitute a derivation from first principles within a specific modified-gravity theory. Rather, it is a phenomenological ansatz motivated by the need to simultaneously explain (i) the rapid mass assembly observed at $z > 10$, (ii) the ultra-compact nature of LRDs, and (iii) their distinctive V-shaped SEDs. The success of this parameterization in fitting the data should be interpreted as evidence that \textit{some form} of enhanced early gravity (or physically equivalent mechanism) is consistent with the observations, not as proof of any particular underlying theory.

Second, the Gaussian window function adopted in Eq.~(1) represents one of many possible choices for the temporal profile of $G_{\rm eff}$. While the Gaussian form provides a natural interpretation in terms of a finite-duration ``enhanced-gravity epoch'' centered on cosmic dawn, we cannot rule out alternative profiles that might fit the data equally well. Future work should explore the sensitivity of our conclusions to the assumed functional form, ideally within a Bayesian model comparison framework that marginalizes over a family of possible profiles.

Finally, we note that the relic hypothesis makes additional predictions beyond those tested here---most notably, the expectation that genuine relic objects should exhibit characteristically high stellar velocity dispersions ($\sigma_* \gtrsim 300-500\rm\,km\,s^{-1}$) due to their deep potential wells, and potentially anomalous gravitational lensing signals if sufficiently overdense. Upcoming JWST/NIRSpec medium-resolution spectroscopy of representative LRDs will provide critical tests of these ancillary predictions.

\section{Conclusion} \label{sec:conclusion}

Through a rigorous multi-wavelength statistical evaluation of 260 confirmed Little Red Dots, we have resolved the crippling radiation paradox imposed by their absent X-ray and mid-infrared features. Dispensing with unphysical dust models, we demonstrate that a time-dependent early gravitational enhancement ($G_{\rm eff}$) results in stabilized, ultra-compact relics. The proposed Gravitational Relic Model provides the required phenomenological spectral displacement ($z_{\rm dist} = 0.1700$), comprehensively erasing the missing mid-infrared budget and effortlessly translating into broad Keplerian linewidths without simultaneous X-ray production. Corroborated by the latest detections of extraordinary AGN outflows \citep{Korber2026} and highly overmassive metal-poor seeds \citep{Ivey2026}, LRDs represent not merely peculiar galaxies, but fundamentally frozen artifacts of dynamic structural assembly in the $z>10$ universe.

\section{Data Availability and Reproducibility} \label{sec:data_availability}

To ensure complete statistical transparency and reproducibility, all raw output matrices and supplementary visualizations underpinning this study are made publicly available. 
The complete 260-source parameter scanning results (\texttt{Bulletproof\_Results.csv}) and the stratified empirical Kolmogorov-Smirnov parameter distribution visualizations (\texttt{Three\_Group\_Distribution\_Analysis.pdf}) are hosted on Zenodo [DOI: 10.5281/zenodo.19571370].
The phenomenological grid scanning routines, including SED reconstruction algorithms and multi-dimensional statistical plotting scripts (as demonstrated in the supplemental materials), have been containerized and are available via a dedicated GitHub repository [Link: \url{https://github.com/summicron352-tech/LRD-Relic-Gravity-Fitting}]. 
Researchers are explicitly encouraged to manipulate the generic $G_{\rm eff}$ constraints and directly test alternative parametric correlations within the provided dataset.

\bibliographystyle{aasjournal}
\bibliography{LRDs_Paper_Draft_v2}

\appendix

\section{Full Sample Scan Methodology} \label{sec:appendix_a}

To systematically discriminate between the standard dust attenuation framework and the Gravitational Relic Hypothesis, we compiled a parent sample of 260 LRDs \citep{Kokorev2024} (hereafter \texttt{Kokorev\_LRDs\_Full}). The sample strictly mandates existing point-like morphological constraints from JWST/NIRCam and spectroscopic coverage from JWST/NIRSpec. We evaluated each source using an independent comprehensive 5-parameter spectral energy distribution (SED) grid scanning routine, informed by standard Bayesian inference frameworks \citep{EAZY2010, Prospector2019, BAGPIPES2020}.

The two competing paradigms were parameterized identically across shared components (e.g., stellar mass, unattenuated power-law slope $\alpha_{\rm AGN}$) to ensure statistical purity. The fundamental bifurcation occurs in treating the observed V-shaped red break:
1. \textbf{The Null Model} adheres to standard $\Lambda$CDM astrophysics. It assumes $z_{\rm dist} = 0$ (no gravitational redshift distortion) and forces the optimizer to fit the extremely red optical slopes utilizing purely physical dust attenuation arrays (characterized by $A_V$, e.g., \citealt{Calzetti2000}).
2. \textbf{The Relic Model ($G_{\rm eff}$)} leverages the theoretical displacement $z_{\rm dist} = 0.1700$ independently derived from our effective gravity window framework. Dust ($A_V$) in this model is allowed to relax to natural ambient physical lower bounds.

For each of the 260 sources, the Python-based fitting architecture explored millions of template permutations, extracting the absolute $\chi^2$ minimum for both models. To account for varying model dimensionalities, we relied on the Reduced Chi-Squared ($\chi^2_\nu$) and the Bayesian Information Criterion (BIC), aggressively penalizing the excessive dust complexity enforced by the Null Model. A statistical ``verdict'' was explicitly assigned based on the magnitude of the $\Delta \chi^2_{\nu}$ dominance. 

\begin{figure}[ht!]
\centering
\includegraphics[width=\linewidth]{ParamScan_Heatmap_AGamma.png}
\caption{Convergence topography of the dense 5-dimensional grid search evaluating the entire empirical LRD base. Absolute convergence securely anchors our adopted parameters across the massive array space minimizing degenerate trapping points.}
\label{fig:param_heatmap}
\end{figure}

\section{Complete Scan Results for 260 LRDs} \label{sec:appendix_b}

The complete machine-readable table detailing the 260 source scan results (\texttt{Bulletproof\_Results.csv}) is provided as supplementary material attached to this manuscript. 

To summarize the macroscopic distribution:
\begin{itemize}
    \item \textbf{Support (76 sources, 29.2\%):} Distinct and overwhelming statistical preference for the Relic Model ($\Delta \chi^2_{\nu} \gg 10$). These sources demand heavily unphysical $A_V$ parameters under the Null Model, which completely vanish under the introduction of $z_{\rm dist}$.
    \item \textbf{Weak Support (35 sources, 13.5\%):} Consistent statistical preference for the Relic Model.
    \item \textbf{Neutral (106 sources, 40.8\%):} Both models yield statistically indistinguishable residuals. These sources generally have moderately red slopes that do not strictly violate conventional dust boundaries, making both assumptions mathematically viable. According to BIC, these inherently ``default'' to the Relic model due to lower parametric tension.
    \item \textbf{Null Wins (43 sources, 16.5\%):} These exceptional sources explicitly favor the standard dust model. Detailed SED profile analyses suggest that these objects represent LRD-like interloper galaxies at strictly lower genuine redshifts ($z < 4$), or genuinely diverse dusty starbursts entirely unrelated to the high-$G$ formation era.
\end{itemize}

Collectively, the $G_{\rm eff}$ relic framework provides an acceptable or preferred description for 83.5\% (217/260) of the complete empirical sample. Of these, 42.7\% show statistically significant preference for the relic model, representing robust individual detections. The additional 40.8\% of sources---while not individually discriminating---are consistent with the relic framework and contribute to the aggregate statistical signal when analyzed as an ensemble. Only 16.5\% of the sample is definitively better described by standard dust extinction, establishing that the relic hypothesis cannot be dismissed as a universal explanation for the LRD population.

\section{Grid SED Fitting Gallery} \label{sec:appendix_c}

To comprehensively satisfy verification prerequisites typically enforced within statistical phenomenological investigations, we visibly plot the individual model convergences. Figure \ref{fig:appendix_sed_grid} constitutes a comprehensive visual mosaic highlighting randomly selected subsets extracted strictly from the Strong Relic Support paradigm demonstrating the catastrophic divergence of typical $\Lambda$CDM assumptions against high-fidelity JWST constraints when applied to Early Universe artifacts.

\begin{figure*}[ht!]
\centering
\includegraphics[width=\textwidth]{Appendix_Grid_SED_Final.png}
\caption{Comprehensive Grid SED atlas delineating extensive sub-sample convergences strictly under the Gravitational Relic formalism versus the chromatic dust Null configuration.}
\label{fig:appendix_sed_grid}
\end{figure*}

\begin{figure*}[ht!]
\centering
\includegraphics[width=0.32\textwidth]{LRD_Visualization_PNGs/page-5.png}
\hfill
\includegraphics[width=0.32\textwidth]{LRD_Visualization_PNGs/page-6.png}
\hfill
\includegraphics[width=0.32\textwidth]{LRD_Visualization_PNGs/page-7.png}
\caption{Complete Supplemental Diagnostics (Part V, VI, \& VII). Outlining edge cases, failure regimes (e.g., standard starburst interlopers classified under 'Null Wins'), and extended multidimensional parameter density distributions mapping the entirety of the 260 empirical JWST observations.}
\label{fig:viz_pages_5_6_7}
\end{figure*}

\end{document}
